% !TEX program = xelatex
% Compile twice with XeLaTeX so that references are resolved correctly.
\documentclass[conference]{IEEEtran}

\usepackage{amsmath}
\usepackage{booktabs}
\usepackage{graphicx}
\usepackage{kotex}
\usepackage{tabularx}
\usepackage[hidelinks]{hyperref}
\graphicspath{{figures/}{docs/figures/}}

\title{HybridMLQ: 전용 큐 보호와 작업 탈취를 결합한\\
비선점 체크인 스케줄링}
\author{}

\begin{document}

\maketitle
\suppressfloats[t]

\begin{abstract}
등급별 전용·공용 공항 체크인 카운터에서는 전용 큐 우선 원칙과 전체 평균
반환 시간(ATT) 사이의 상충을 고려해야 한다. 본 연구는 비선점 Priority,
Shortest Job First (SJF), Highest Response Ratio Next (HRRN)를 결합한
온라인 휴리스틱 HybridMLQ를 제안한다. 과제의 승객 50명 고정 입력에서 ATT는
15.88로 FCFS의 19.52와 Fixed-Priority의 22.50보다 낮았지만, 전역 비선점
SJF의 14.90보다는 6.58\% 높았다. 반면 최악 등급 ATT는 SJF의 44.25에서
31.38로 낮아졌다. 따라서 HybridMLQ는 전체 ATT 최적화보다 전용 큐 보호와
등급 간 상충 조정을 택한 정책이며, 일반적 우월성을 주장하지 않는다.
\end{abstract}

\begin{IEEEkeywords}
비선점 스케줄링, Multi-Level Queue, SJF, HRRN, 작업 탈취
\end{IEEEkeywords}

\section{문제 정의}

전용 카운터를 엄격히 운영하면 자기 등급 큐가 비었을 때 자원이 유휴 상태가
되고, 모든 승객을 통합하면 등급별 서비스 정책이 사라진다. 본 연구의 목표는
전용 큐 우선과 작업 탈취를 결합한 정책을 설계하고, 지정된 세 baseline 대비
전체 ATT와 등급별 ATT의 상충을 분석하는 것이다. 스케줄링 용어는
Pinedo~\cite{pinedo}를 따른다.

시스템은 동일 속도의 카운터 5개(C1: FIRST, C2: BUSINESS, C3: ECONOMY,
C4--C5: FLEX)와 승객 50명으로 구성된다. 도착 시각은 0--58이고 등급별
인원은 8, 10, 32명이며, 평균 서비스 시간은 각각 12.8, 9.1, 5.7이다.
승객 $i$의 도착 시각 $a_i$, 서비스 시간 $s_i$, 완료 시각 $C_i$에 대해

\begin{equation}
T_i=C_i-a_i, \qquad
\mathrm{ATT}=\frac{1}{n}\sum_{i=1}^{n}(C_i-a_i)
\label{eq:att}
\end{equation}

로 정의한다. 서비스는 비선점형이고, 모든 선택은 현재까지 도착한 승객만
사용한다. 여기서 전용 큐 ``보호''는 자기 등급 승객이 기다리는 동안 다른
등급을 선택하지 않는다는 규칙이며, 등급별 ATT 보장을 뜻하지 않는다.

\section{HybridMLQ}

도착 승객은 FIRST, BUSINESS, ECONOMY 큐에 저장된다. 빈 카운터는
표~\ref{tab:dispatch}의 규칙으로 선택한다. 비선점 Priority는 전용 큐를
우선하고, SJF는 짧은 작업을 선택해 convoy effect를 줄이며, 제한적 HRRN은
긴 Economy 작업의 starvation을 완화한다. 즉 세 알고리즘은 각각 큐 간
정책, 큐 내부 효율, aging을 담당한다.

\begin{table}[t]
\caption{HybridMLQ 디스패치 규칙}
\label{tab:dispatch}
\centering
\footnotesize
\begin{tabularx}{\columnwidth}{@{}lX@{}}
\toprule
카운터 & 선택 규칙 \\
\midrule
C1 & FIRST에서 SJF; 비면 나머지 큐에서 SJF \\
C2 & BUSINESS에서 SJF; 비면 나머지 큐에서 SJF \\
C3 & ECONOMY에서 SJF/aging; 비면 나머지 큐에서 SJF \\
C4, C5 & 전체 큐에서 SJF \\
\bottomrule
\end{tabularx}
\end{table}

C3는 10 tick 이상 기다린 Economy 승객이 있을 때만 그 후보 중 HRRN을
적용하며, 그렇지 않으면 SJF를 사용한다. 시각 \textit{now}에서

\begin{equation}
W_i=\mathit{now}-a_i, \qquad RR_i=\frac{W_i+s_i}{s_i}.
\label{eq:hrrn}
\end{equation}

SJF 동률은 등급, 도착 시각, 승객 ID 순으로, HRRN 동률은 서비스 시간과
ID 순으로 해소한다~\cite{implementation}. 그림~\ref{fig:gantt}의 P01과
P40은 전용 큐가 비었을 때 발생한 두 작업 탈취 사례이다.

\begin{figure}[t]
\centering
\includegraphics[width=\columnwidth]{hybridmlq_gantt.pdf}
\caption{고정 입력에서 HybridMLQ가 생성한 실행 일정}
\label{fig:gantt}
\end{figure}

\section{구현 및 평가}

시뮬레이터는 각 tick에 도착, 완료, C1--C5 순 디스패치, 잔여 시간 감소를
수행하고 상태 변화를 기록한다. 스케줄러는 교체 가능한 Strategy 모듈이며,
$C_i-\mathit{start}_i=s_i$를 검사해 비선점을 검증한다~\cite{implementation}.
과제 정의에 따라 FCFS는 전역 도착 순, Fixed-Priority는 전역
FIRST--BUSINESS--ECONOMY 및 등급 내 FCFS, SJF는 등급 무관 전역 최단
서비스 시간 순으로 구현했다. 네 정책은 동일 입력과 엔진에서 비교했다.

\begin{table}[t]
\caption{알고리즘별 평균 반환 시간}
\label{tab:results}
\centering
\resizebox{\columnwidth}{!}{%
\begin{tabular}{@{}lrrrr@{}}
\toprule
알고리즘 & 전체 ATT & FIRST & BUSINESS & ECONOMY \\
\midrule
FCFS                  & 19.52 & 23.88 & 20.36 & 18.10 \\
Fixed-Priority        & 22.50 & \textbf{13.25} & \textbf{11.09} & 28.94 \\
Non-preemptive SJF    & \textbf{14.90} & 44.25 & 16.18 & \textbf{6.87} \\
HybridMLQ             & 15.88 & 31.38 & 24.18 & 8.94 \\
\bottomrule
\end{tabular}%
}
\end{table}

HybridMLQ는 FCFS와 Fixed-Priority보다 전체 ATT를 각각 18.65\%, 29.42\%
줄였다. 이는 짧은 Economy 작업이 64\%인 입력에서 SJF가 FCFS의 convoy
effect와 Priority의 Economy 지연을 줄이는 경향과 일치한다. 그러나 전역
SJF는 등급 보호 없이 짧은 작업을 선택하므로 ATT가 14.90으로 가장 낮았고,
평균 서비스 시간이 가장 긴 FIRST의 ATT는 44.25까지 증가했다. HybridMLQ는
전체 ATT를 0.98 희생하는 대신 최악 등급 ATT를 31.38로 낮췄다.

구성 요소 제거 시 ATT는 작업 탈취 없이 16.18, aging 없이 16.02였다.
Aging은 Economy ATT를 10.61에서 8.94로 낮추는 대신 BUSINESS ATT를
19.91에서 24.18로 높였다. 이 고정 입력에서는 두 요소의 기여와 대기 시간
이동이 관찰되지만 일반화할 수는 없다. 그림~\ref{fig:tradeoff}는 전체 ATT와
최악 등급 ATT의 상충을 요약한다.

\begin{figure}[t]
\centering
\includegraphics[width=\columnwidth]{performance_tradeoff.pdf}
\caption{전체 ATT와 최악 등급 ATT의 상충 관계}
\label{fig:tradeoff}
\end{figure}

\section{한계 및 결론}

평가는 단일 데이터셋과 간단한 제거 실험에 한정되며, aging 임계값 민감도와
통계적 반복 검증이 없다. 최적해나 하한이 없어 optimality gap을 알 수 없고,
고정된 카운터 순서, 비선점, 동일 처리 속도도 실제 환경을 단순화한다.

HybridMLQ는 과제의 세 비선점 알고리즘 조합과 카운터 배정 전략을 구현해
FCFS와 Fixed-Priority보다 낮은 ATT를 기록했다. 전역 SJF보다는 느리지만
최악 등급 ATT를 완화했으며, 이는 전체 효율과 전용 큐 보호 사이의 명시적인
trade-off이다. 다양한 부하와 임계값에서 재검증해야 일반화할 수 있다.

\begin{thebibliography}{00}

\bibitem{pinedo}
M. L. Pinedo, \emph{Scheduling: Theory, Algorithms, and Systems}.
Springer.

\bibitem{implementation}
SkyPort, ``SkyPort simulator implementation and tests,'' source code.

\end{thebibliography}

\end{document}
